\section{Introduction}

A succinct non-interactive argument of knowledge (SNARK) allows a
prover to certify a computation with a short proof that can be
verified substantially faster than re-executing the computation.
Zero-knowledge variants additionally hide the witness.
These properties make SNARKs useful for verifying computations
performed by an untrusted party~\cite{eurocrypt/Groth16,crypto/Setty20}.
The cost of generating the proof, however, can substantially exceed
the cost of the computation being certified.

Soundness imposes an additional constraint on the arithmetic used
by the proof.
Many efficient SNARKs use a polynomial interactive oracle proof
(PIOP) to reduce correctness to polynomial identities tested at random
field points~\cite{crypto/Setty20,eurocrypt/ChenBBZ23}.
The probability that a false identity passes such a test is controlled
by its degree relative to the size of the challenge field.
To make this probability negligible, standard constructions use
challenges from a large prime field or a sufficiently large extension
of the witness field.
The prover must then perform arithmetic involving these challenges,
even when the original witness consists only of short integers.

Embedding a short integer into a large field preserves its value,
but does not preserve the cost of computing with it throughout the
proof. In a conventional PIOP, challenge-dependent linear combinations
generally produce full-field values.
Subsequent multiplications therefore require large-field arithmetic,
whose cost reflects the field representation rather than the initial
witness bit width~\cite{iacr/BagadDDT25}.
Arithmetic modulo a multiword prime, or in a high-degree extension,
can be much more expensive than the native short-integer
multiplications being certified.
Materializing the resulting full-field vectors also expands the
prover's working memory.

Quantized machine-learning inference makes this mismatch concrete.
Low-precision inference uses short integers to reduce computation
and memory costs: SmoothQuant, for example, supports matrix
multiplications with 8-bit weights and 8-bit
activations~\cite{icml/XiaoLSWDH23}.
These representations also appear in zero-knowledge machine learning
(zkML). Kang et al.~\cite{corr/KangHSS22} prove DNN inference with
8-bit weights and activations, with intermediate computations of up
to 32 bits, using an arithmetization over a 254-bit prime field.
Thus the precision sufficient for inference can be far below the
field size used by the proof system.
Quantization reduces the cost of executing the model, but the prover
does not automatically inherit the same savings when checking the
relations among its weights and activations.

Exploiting this gap requires more than a compact witness commitment.
A packed polynomial commitment scheme (PCS) can reduce the cost of
committing to short values, while the PIOP still processes long vectors
of full-field elements~\cite{eurocrypt/DiamondP26,iacr/BagadDDT25}.
The central problem is therefore to separate the field used to
represent and process the witness from the field used for soundness:
\begin{quote}
\emph{Can we retain small-field arithmetic for most of the prover's
work while using a large field for soundness?}
\end{quote}

We study this question for vector relations over a small finite
field \(\Field\).
Our approach is to reduce a length-\(N\) relation over \(\Field\)
to a length-\(m\) relation over an extension \(\ExtensionField\),
where \(N=m\kappa\) and \([\ExtensionField:\Field]=\kappa\).
Polynomial packing keeps the row-polynomial multiplications over
\(\Field\) and confines the subsequent extension-field proof to
the shorter vectors.
A packed opening protocol authenticates the resulting
challenge-dependent evaluations against the original witness
commitment. Section~\ref{sec:overview} explains this reduction.

The working field need not be prime. For binary computations,
the tower \(\mathbb F_2\subset\mathbb K=\mathbb F_{2^8}
\subset\ExtensionField\) allows the PIOP to use byte-sized
working-field arithmetic while retaining the larger field for
soundness. We instantiate the reductions and their commitment
interface over \(\mathbb K\), with extension degree
\([\ExtensionField:\mathbb K]\), and retain Booleanity constraints
for binary witness values. This field choice and its parameter
requirements are specified in Section~\ref{sec:binary-packing}.

\subsection{Related Work}
\label{sec:related-work}

\paragraph{SNARKs and polynomial IOPs.}
A common construction separates algebraic verification
from cryptographic authentication: a polynomial interactive oracle
proof (PIOP) reduces the relation to polynomial evaluation claims,
and a polynomial commitment scheme (PCS) authenticates those claims.
For multilinear representations, Spartan~\cite{crypto/Setty20} uses
sumcheck to verify rank-one constraint systems, while
HyperPlonk~\cite{eurocrypt/ChenBBZ23} expresses PLONK-style gate and
permutation constraints over the Boolean hypercube.
The PIOP and PCS contribute different prover costs.
Two prominent cryptographic approaches, elliptic-curve-based and
hash-based, offer different commitment tradeoffs, but the cost of
processing short-valued witnesses also depends on the PIOP.

\paragraph{ECC-based constructions.}
Elliptic-curve cryptography (ECC) supports both pairing-based arguments
and discrete-logarithm-based inner-product arguments.
Groth's pairing-based SNARK~\cite{eurocrypt/Groth16} achieves proofs
of three group elements with circuit-specific preprocessing.
KZG~\cite{asiacrypt/KateZG10} provides constant-size commitments and
evaluation proofs, furnishing a polynomial commitment interface for
pairing-based SNARKs.
Sonic~\cite{ccs/MallerBKM19}, Marlin~\cite{eurocrypt/ChiesaHMMVW20},
and PLONK~\cite{iacr/GabizonWC19} support universal and updatable
structured reference strings.
On the discrete-logarithm side, the inner-product argument underlying
Bulletproofs~\cite{sp/BunzBBPWM18} gives logarithmic communication
without a trusted setup, and Hyrax~\cite{sp/WahbyTSTW18} provides
a multilinear commitment with communication--verification tradeoffs.
For these constructions, the polynomial arithmetic takes place in
the scalar field of the chosen group.
Its size is constrained by cryptographic security, rather than by
the range of the witness entries: representing an entry by a short
integer does not make this ambient field smaller.

\paragraph{Hash-based constructions.}
Hash-based arguments replace group-based commitments with
error-correcting encodings and Merkle authentication.
STARKs~\cite{iacr/Ben-SassonBHR18} exemplify this transparent approach,
with FRI~\cite{icalp/Ben-SassonBHR18} supplying Reed--Solomon
proximity testing.
Ligero~\cite{dcc/AmesHIV23} uses interleaved Reed--Solomon codes
to obtain sublinear arguments without a trusted setup.
Brakedown~\cite{crypto/GolovnevLSTW23} obtains linear-time proving
from linear-time encodable codes over sufficiently large finite fields.
BaseFold~\cite{crypto/ZeilbergerCF24} combines foldable codes with
sumcheck to obtain a field-agnostic multilinear PCS.
WHIR~\cite{eurocrypt/ArnonCFY25} uses constrained Reed--Solomon
proximity testing to support polynomial evaluation queries with
fast verification.
These constructions decouple the arithmetic field from a group order,
allowing field choices subject to their coding and soundness conditions.
This flexibility enables small-field witness encodings, but does not
itself change how a PIOP processes those encodings.

\paragraph{The short-integer bottleneck in the PIOP.}
Neither commitment family, by itself, makes the PIOP bit-width-aware.
The issue is the cost of computing with short-valued witnesses,
even when their range constraints have already been enforced.
In conventional sumcheck, random challenges combine entries of the
witness tables into new field elements.
Such combinations need not retain the original short-integer
representation. With extension-field challenges, even inputs in a
small base field generally yield values outside that base field.
The remaining rounds therefore operate on full-field values unless
the prover uses additional structure~\cite{iacr/BagadDDT25}.
An ECC-based PCS fixes the ambient scalar field, whereas a hash-based
PCS permits more field choices. Neither property alone ensures that
the PIOP's arithmetic and working memory scale with the witness bit
width. Taking challenges from a smaller field also changes the
algebraic soundness error and requires a corresponding protocol
analysis. Commitment compression and bit-width-aware relation
checking are thus separate requirements.

\paragraph{Small-field embeddings and Binius.}
Packing several witness entries into one extension-field element
predates Binius. Cascudo and Giunta~\cite{iacr/CascudoG21} use
reverse multiplication-friendly embeddings to verify Boolean R1CS
through larger-field linear and row checks.
Binius, due to Diamond and Posen~\cite{eurocrypt/DiamondP25},
adapts Brakedown to multilinear commitments over binary towers
without coefficient-wise embedding overhead, and develops
corresponding product, permutation, and lookup arguments.
Their subsequent ring-switching construction~\cite{eurocrypt/DiamondP26}
gives a generic reduction from a multilinear PCS over an extension
field to one over its ground field.
Its commitment cost matches that of the underlying PCS on inputs
of the same bit size; its opening reduction adds linear prover
and logarithmic verifier overhead.
The compiler applies to general field extensions. Their concrete
PCS instantiation uses a binary-field adaptation of BaseFold.
Ring-switching addresses the interface between small-field polynomial
evaluations and packed commitments. Exploiting the same representation
in the circuit PIOP requires relation-specific techniques, as in
Binius's binary-tower arguments.

\paragraph{Small-value and small-field sumcheck.}
Dedicated sumcheck algorithms address this PIOP-level cost directly.
Bagad, Domb, and Thaler~\cite{iacr/BagadDT24} optimize the sumcheck
prover when the input polynomials take values in a small subfield,
replacing expensive extension-field multiplications with additional
base-field arithmetic.
Bagad, Dao, Domb, and Thaler~\cite{iacr/BagadDDT25} extend this line
to both short integers and small-subfield values, combine it with
equality-polynomial optimizations, and integrate the resulting prover
into Jolt.
Dao and Thaler~\cite{iacr/DaoT24} exploit binary-tower structure in
zero-check, reducing equality-table computation and generalizing
constraint packing with linearly independent extension-field weights.
Wei et al.~\cite{iacr/WeiWXZDWZ25} give packed sumcheck over small
prime fields using repetition and folding to keep the computation
in the base field, with an application to the vector relations in
TFHE bootstrapping.
These approaches modify the relation proof itself, through
small-value arithmetic, structured packing, or repetition and folding.
Their gains are additional to those obtained by choosing a faster PCS.

\paragraph{Relation to our work.}
Mikros develops a relation-level reduction for arithmetic over small
finite fields, together with the packed opening interface it needs.
For inner- and Hadamard-product relations, we arrange a length-\(N\)
witness as an \(m\times\kappa\) matrix over a small field
\(\Field\), where \(N=m\kappa\) and
\([\ExtensionField:\Field]=\kappa\).
We encode each row by a polynomial over \(\Field\) and express the
relation as a low-degree polynomial identity.
Evaluating this identity at an extension-field challenge leaves a
length-\(m\) inner-product claim over \(\ExtensionField\).
Unlike commitment packing alone, this step compresses the relation
passed to the extension-field proof.
It retains the base-field row multiplications while
restricting the subsequent extension-field proof to the shorter
vectors.
The commitment uses a fixed field basis, whereas the relation proof
requires challenge-dependent packing weights.
Our opening reduction connects the two through base-field coordinate
identities and terminates in an opening of an underlying multilinear
PCS over \(\ExtensionField\).
Sections~\ref{sec:overview} and~\ref{sec:pcs} develop these
relation and opening reductions, respectively.
